Representation similarity methods are widely used to analyze deep networks, but controlled studies on plain multilayer perceptrons (MLPs) remain less common than analogous work on convolutional or attention-based models. We ask a simple question: are hidden states in different MLP layers effectively unrelated, or do they preserve enough shared structure that layer-wise comparison is meaningful? To answer this, we train matched four-hidden-layer ReLU MLPs on \mnist, \fashionmnist, and \cifar using three random seeds and compare hidden activations with linear \cka, \svcca, \pwcca, and sample-wise cosine similarity. We complement these geometric comparisons with linear probes on every hidden layer.

Our results show a clear within-model locality pattern. Adjacent layers are more similar than non-adjacent layers in all 12 dataset-metric comparisons, with 10 reaching $p<0.05$. Cross-seed comparisons are weaker: same-depth layers exceed off-depth layers in 8 of 12 cases, and only 5 are statistically significant. Linear probes show that later layers usually improve decodable task information, from $0.8767$ to $0.9108$ on \mnist and from $0.8308$ to $0.8450$ on \fashionmnist, while \cifar peaks at layer 3 with $0.4292$.

These findings argue against the view that MLP hidden spaces are completely rewritten at each layer. Instead, plain feedforward MLPs appear to transform representations progressively, although the strength of this conclusion depends on which similarity metric is used.
